\section{Key Decisions and Assumptions}

\subsection{Missing Review Information}

The original dataset contained 10,052 missing values in both
\texttt{last\_review} and \texttt{reviews\_per\_month}. Inspection showed that
these missing values were associated with listings that had zero recorded
reviews.

Missing \texttt{reviews\_per\_month} values were therefore replaced with zero
only where \texttt{number\_of\_reviews} was also zero. This preserved
unreviewed listings while representing their observed monthly review activity
numerically. Missing \texttt{last\_review} values were retained because a
listing with no reviews has no meaningful last-review date.

This variable-specific approach was preferable to deleting every incomplete
record. Missing-data treatment can affect representativeness and introduce
bias, so the meaning and pattern of missingness should be considered before
selecting a treatment \parencite{kang2013}.

The small numbers of missing listing names and host names were retained because
those descriptive fields were not needed for the principal price and
availability analyses.

\subsection{Invalid and Extreme Values}

The cleaning process removed 11 records whose prices were zero. A non-positive
nightly price was treated as invalid for the planned price analysis. This
reduced the dataset from 48,895 to 48,884 listings.

Large positive prices and long minimum-stay requirements were not automatically
removed. Although these observations are unusual, the dataset provides
insufficient evidence to classify them as errors. Removing them without an
explicit validity rule could exclude legitimate listings and alter the observed
distribution.

The price distribution was strongly right-skewed. The full cleaned dataset had
a mean price of \$152.76 and a median price of \$106.00. For listings priced at
or below the 99th percentile of \$799, the mean decreased to \$137.58, whereas
the median changed only to \$105.00. Median prices were consequently used for
the main grouped comparisons because they were less sensitive to extreme
values.

Prices above the 99th percentile were excluded only from Figure 2 in the accompanying notebook
to improve its readability. They remained in the cleaned dataset and in the
other analyses.

\subsection{Choice of Exploratory Comparisons}

Neighbourhood group and room type were selected as the central categorical
variables because they allowed direct comparison of location and accommodation
category. Price and availability were grouped by both variables so that their
combined patterns would not be hidden by separate aggregate summaries.

The correlation matrix was included to examine linear relationships among the
numerical variables. Correlation coefficients were interpreted as descriptive
associations and not as evidence of causation.